% Converted from README.md
\subsection{tmux fuzz testiranje}\label{tmux-fuzz-testiranje}

Projekat \texttt{tmux} je dobar primer primene fuzz testiranja nad programom koji obrađuje dosta tekstualnih protokola i terminalskih sekvenci. Fuzz testovi su smešteni u direktorijumu \href{https://github.com/tmux/tmux/tree/master/fuzz}{\texttt{fuzz}} u uzvodnom repozitorijumu.

Ovi testovi su korisni za traženje grešaka u parserima i kodu koji obrađuje ulaz iz terminala, konfiguracionih datoteka i mrežnih veza. Takav kod često ima veliki broj grana i mnogo kratkih, nevalidnih ulaza koji se retko pokriju ručno pisanim testovima.

Prilikom čitanja ovog primera obratiti pažnju na tri stvari:

\begin{itemize}
\tightlist
\item
  fuzz target izoluje deo programa koji obrađuje jedan format ulaza;
\item
  ulaz fuzzera se tretira kao niz bajtova, bez pretpostavke da je sintaksno ispravan;
\item
  pad programa, prekoračenje memorije ili beskonačno izvršavanje predstavljaju koristan nalaz.
\end{itemize}

Ako se ovakav pristup prenosi na sopstveni projekat, prvo treba izabrati malu funkciju koja parsira ulaz i napraviti omotač koji samo priprema stanje, poziva parser i oslobađa resurse. Fuzzer tada može da izvrši veliki broj kratkih testova bez pokretanja cele aplikacije.
